\documentclass{beamer}
\usetheme{Madrid}

% ---- latex_handler theme support ------------------------------------------
% Beamer has its own colour themes, so switch those in the preamble as well.
\providecommand{\lhtheme}{light}
\def\lhdarkvalue{dark}
\ifx\lhtheme\lhdarkvalue
  \usecolortheme{dolphin}
\else
  \usecolortheme{seahorse}
\fi
\definecolor{lhlightaccent}{HTML}{1F4E79}
\definecolor{lhdarkaccent}{HTML}{E8A33D}
\colorlet{lhaccent}{lhlightaccent}
\def\lhthemevalue{Light}
\ifx\lhtheme\lhdarkvalue
  \colorlet{lhaccent}{lhdarkaccent}
  \def\lhthemevalue{Dark}
\fi
\providecommand{\lhbar}{}
\providecommand{\lhfootertext}{}
\InputIfFileExists{lh-branding}{}{}
% ---------------------------------------------------------------------------

\title{Presentation Variation}
\subtitle{Built from \texttt{latex-templates}}
\author{Your Name}
\date{\today}

\begin{document}

\frame{\titlepage}

\begin{frame}{Outline}
  \begin{itemize}
    \item Why a multi-document template is convenient
    \item How \texttt{latex\_handler} discovers documents
    \item How the light and dark themes differ
  \end{itemize}
\end{frame}

\begin{frame}[fragile]{Discovery}
  A file named \texttt{main.tex} or \texttt{*\_main.tex} under \texttt{src/} is
  compiled automatically. Compile this deck with:
  \begin{verbatim}
    latex_handler compile --file src/beamer_main.tex
  \end{verbatim}
\end{frame}

\begin{frame}{Theme}
  The active theme is \alert{\lhthemevalue}. The \texttt{dolphin} colour theme is
  used for dark builds and \texttt{seahorse} for light builds.
  \lhbar
\end{frame}

\begin{frame}{Summary}
  \lhfootertext
\end{frame}

\end{document}
